\documentclass{article}
\usepackage[utf8]{inputenc}
\usepackage{amsmath}
\usepackage{geometry}
\geometry{margin=2cm}
\begin{document}

\section*{Questão 95 — ENEM 2024 — Ciências da Natureza}

A tirinha apresenta esquimós dentro de um iglu, habitação hemisférica comum nas regiões setentrionais da Groenlândia, Canadá e Alasca. Apesar do ambiente externo extremamente frio, a água não congela dentro do iglu; por isso, uma geladeira é utilizada para fazer gelo.

\section*{Interpretação e Estratégia}

A questão pede para justificar por que a geladeira é necessária para fabricar gelo dentro do iglu, mesmo sendo um ambiente naturalmente frio. Deve-se analisar a temperatura interna do iglu e suas propriedades térmicas para entender a mudança de estado da água.

\subsection*{Análise}
\begin{itemize}
  \item O iglu funciona como isolante térmico, mantendo a temperatura interna acima de 0~\degree C para conforto.
  \item A água congela a 0~\degree C ou menos sob pressão atmosférica normal.
  \item Logo, a água dentro do iglu não congela espontaneamente e necessita de equipamento que a resfrie abaixo de zero.
\end{itemize}

\section*{Conteúdo e Mini-Revisão}

\begin{table}[h!]
\centering
\begin{tabular}{|p{5cm}|p{5cm}|p{5cm}|}
\hline
\textbf{Conceito} & \textbf{Ideia Chave} & \textbf{Aplicação} \\
\hline
Ponto de Congelamento da Água & Água congela a 0~\degree C & Temperatura interna do iglu acima de 0~\degree C mantém água líquida \\
\hline
Isolamento Térmico do Iglu & Mantém temperatura estável confortável & Ambiente interno protegido do frio extremo externo \\
\hline
Mudança de Fase (Líquido para Sólido) & Depende principalmente da temperatura & Necessário resfriar água a menos de 0~\degree C para gelo \\
\hline
\end{tabular}
\caption{Revisão dos conceitos físicos relevantes}
\end{table}

\section*{Resolução e "Pulo do Gato"}

A geladeira é necessária porque o iglu mantém uma temperatura interna um pouco acima de zero grau Celsius, o que impede que a água congele naturalmente. Somente com o uso da geladeira a água pode ser resfriada abaixo do ponto de congelamento, originando o gelo utilizado pelas crianças.

\textbf{Pulo do gato:} a temperatura interna, não o frio externo, é que determina a necessidade da geladeira para congelar a água dentro do iglu.

\section*{Armadilhas e Conceitos Equivocados}

\begin{itemize}
  \item Considerar que a umidade do ambiente impede a formação de gelo.
  \item Confundir isolamento térmico com impedimento total do resfriamento.
  \item Atribuir a ausência de gelo à circulação de ar ou convecção.
  \item Pensar que a pressão dentro do iglu altere significativamente o ponto de congelamento.
\end{itemize}

\section*{Código da Questão}

\texttt{CN\_MUDANCAFÍSICA\_MULTIPLA\_001}

\bigskip

\begin{center}
\textbf{Fim da Resolução}
\end{center}


\end{document}